\documentclass[11pt,a4paper]{article}
\usepackage{curso}
\begin{document}
\portada{Lección}{8.1}{Geometría de la información y gradiente natural}

\noindent\textbf{Unidad 8.}\quad\textbf{Objetivos:} la información de Fisher como métrica local de la KL; el gradiente natural como descenso más empinado en
el espacio de distribuciones; comparación con el descenso de gradiente.
\textbf{Prerrequisitos:} 0.3, 6.1.

\section{Fisher como métrica}
Para $p_\theta$ y un paso $d$:
\[\mathrm{KL}(p_{\theta+d}\|p_\theta)=\tfrac12d^\top Fd+O(\|d\|^3),\qquad F=\mathbb E\big[\nabla\log p\,\nabla\log p^\top\big].\]
$F$ mide distancias entre modelos por sus predicciones, no por sus parámetros.

\section{Gradiente natural}
Minimizar $\ell(\theta+d)$ sujeto a $\frac12d^\top Fd=\epsilon$ (Lagrange) da $d\propto-F^{-1}\nabla\ell$:
\[\tilde g=F^{-1}\nabla\ell,\qquad\theta_{k+1}=\theta_k-\eta F^{-1}\nabla\ell.\]
Para pasos pequeños es invariante ante reparametrizaciones. Fisher empírico: $\bar F=JJ^\top$; Gauss--Newton generalizado: $G=J_fH_LJ_f^\top$ (8.2). Adam usa
solo la diagonal del Fisher empírico.

\section{Ejemplo gaussiano \evS}
Pérdida poblacional de $\N(\mu,\sigma^2)$ con datos de media $m=2$ y desviación $s=0{,}1$:
$\ell(\mu,\sigma)=\log\sigma+\frac{(\mu-m)^2+s^2}{2\sigma^2}$, $F=\mathrm{diag}(1/\sigma^2,\,2/\sigma^2)$, de modo que
$F^{-1}\nabla\ell=(\sigma^2\partial_\mu\ell,\ \tfrac{\sigma^2}2\partial_\sigma\ell)$. Desde $(0,1)$ hasta $\ell-\ell^\ast<10^{-3}$: descenso de gradiente con la mejor tasa
fija (barrido de 51 valores): 106 pasos ($\eta\approx0{,}032$); gradiente natural con $\eta=1$: 7 pasos. NGD corrige $\mu$ en un paso (el primer paso da $\mu=m$
exactamente) y luego ajusta $\sigma$.

\section{El costo}
$F\in\R^{N\times N}$ con $N\sim10^6$--$10^9$: formarla cuesta $O(N^2)$ y resolver $O(N^3)$. Las lecciones 8.2--8.5 lo hacen asequible (CG, Woodbury, TNGD, K-FAC).

\begin{aplicacion}
Entrenamiento de redes neuronales; gradiente natural de políticas (TRPO) en aprendizaje por refuerzo; inferencia variacional con gradiente natural.
\end{aplicacion}

\begin{error}
\begin{itemize}[nosep]
\item Confundir Fisher y hessiano. \quad Olvidar el amortiguamiento ($F$ singular si $b\ll N$).
\item Esperar invariancia exacta con pasos finitos. \quad Creer que Adam es un método de gradiente natural.
\end{itemize}
\end{error}

\begin{thebibliography}{9}\small
\bibitem[36]{amari} S.-I. Amari, Natural gradient works efficiently in learning, \emph{Neural Computation} 10, 251 (1998).
\bibitem[37]{martens} J. Martens, New insights and perspectives on the natural gradient method, \emph{JMLR} 21 (2020).
\bibitem[40]{tngd} K. Donatella et al., Thermodynamic natural gradient descent, 2024.
\end{thebibliography}
\end{document}
